% DOCUMENT CLASS SETUP %
\documentclass[12pt, a4paper]{article}

% DOCUMENT PACKAGES %
\usepackage[a4paper, margin=1in]{geometry}
\usepackage{titling}
\usepackage{enumitem}
\usepackage{amsmath}
\usepackage{amssymb}
\usepackage{mathtools}
\usepackage{graphicx}
\usepackage{xcolor}
\usepackage{listings}
\usepackage{hyperref}
\usepackage{csquotes}
\usepackage{tocloft}
\usepackage{array}

% FONT AND ENCODING SETUP %
\usepackage{fontspec}
\usepackage{polyglossia}
\usepackage[Thai]{ucharclasses}
\setdefaultlanguage{english}
\setotherlanguage{thai}
\setTransitionsFor{Thai}{\thaifont}{\normalfont}
\newfontfamily{\thaifont}{THSarabun.ttf}[
    Path=./fonts/,
    BoldFont=THSarabunBold.ttf,
    ItalicFont=THSarabunItalic.ttf,
    BoldItalicFont=THSarabunBoldItalic.ttf,
    Scale=1.23
]
\XeTeXlinebreaklocale "th"
\XeTeXlinebreakskip = 0pt plus 1pt

% CODE BLOCK PACKAGES %
\usepackage{tcolorbox}
\tcbuselibrary{minted, skins, breakable}

% DEFINE GITHUB DARK THEME COLORS %
\definecolor{ghback}{HTML}{0d1117}
\definecolor{ghtext}{HTML}{c9d1d9}
\definecolor{ghborder}{HTML}{30363d}

% CODE BLOCK CONFIGURATION %
\renewcommand{\theFancyVerbLine}{\textcolor{white!60!gray}{\scriptsize\arabic{FancyVerbLine}}}
\newtcblisting{codeblock}[3][]{
    enhanced,
    breakable,
    colback=ghback,
    colframe=ghborder,
    coltitle=ghtext,
    colupper=ghtext,
    fonttitle=\bfseries\small,
    arc=3mm,
    boxrule=0.5mm,
    % PADDING SETTINGS %
    left=25pt,
    right=10pt,
    top=5pt,
    bottom=5pt,
    % MINTED SETTINGS %
    listing engine=minted,
    listing only,
    minted language=#2,
    minted options={
        linenos,
        numbersep=5pt,
        fontsize=\footnotesize,
        breaklines=true,
        breakanywhere=true,
        autogobble,
        style=monokai,
    },
    title={#3},
    #1
}

% DOCUMENT TITLE SETUP %
\title {
	\textbf{2110999 INTRODUCTION TO \LaTeX} \\
    \vspace{1em}
	\Large{\textbf{\LaTeX\ Document Template}} \\
    \vspace{0.5em}
    \large{Computer Engineering, Chulalongkorn University}
}
\author{Worralop Srichainont}
\date{\today}
\setlength{\droptitle}{15em}

% DISABLE PAGE NUMBER IN TABLE OF CONTENTS
\tocloftpagestyle{empty}

% DOCUMENT CONTENTS
\begin{document}

    % COVER PAGE %
    \maketitle
    \thispagestyle{empty}
    \clearpage

    % TABLE OF CONTENTS PAGE %
    \tableofcontents
    \thispagestyle{empty}
    \clearpage

    % DOCUMENT BODY %
    \pagenumbering{arabic}

    \part{Introduction}

        This document serves as Worralop Srichainont's personal \LaTeX\ template for lecture notes and assignment reports.
        You can find the source code published in my \href{https://github.com/reisenx/LATEX-TEMPLATE}{GitHub repository}.
        It features a simple cover page, a table of contents, and code snippet formatting using the \texttt{minted} package.
        It also fully supports ภาษาไทย (Thai language) using the THSarabunPSK font.
        You are free to use and modify this template without any restrictions.

        \vspace{1em}

        Before compiling this document, you will need to install the Pygments library for the \texttt{minted} package.
        You can easily do this in your terminal using the following command:
        \begin{codeblock}{bash}{}
            pip install Pygments
        \end{codeblock}

        \vspace{1em}

        For your \LaTeX\ editor, \href{https://www.overleaf.com/}{\textbf{Overleaf}} is the easiest online option.
        If you prefer to work locally, you can install \LaTeX\ on your computer via \href{https://miktex.org/download}{\textbf{MikTeX}}
        and use a text editor such as \href{https://code.visualstudio.com/}{\textbf{Visual Studio Code}}
        or \href{https://www.xm1math.net/texmaker/}{\textbf{TeXMaker}}.
        Please note that using Visual Studio Code requires an additional plugin, such as the LaTeX Workshop extension by James Yu.

        \vspace{1em}

        To successfully compile this document, please select \textbf{XeLaTeX} as the compiler in Overleaf or your local editor's settings.
        Alternatively, you can manually run this command in your terminal:
        \begin{codeblock}{bash}{}
            xelatex -shell-escape -interaction=nonstopmode template.tex
        \end{codeblock}

        \vspace{1em}

        Lastly, the remainder of this document serves as a practical demonstration of various \LaTeX\ commands.
        It includes examples of text formatting, lists, tables, mathematical notation, images, and code snippets.
        Feel free to explore the following pages as a helpful reference.
        I hope you find this template useful for your coursework and projects, and please enjoy using it!

    \pagebreak

    \part{Text Formatting}
        \section{Font Styles}
            \LaTeX\ has various of font styles just like other document editor such as Microsoft Word.

            \subsection{English}

                \begin{table}[h]
                    \centering
                    \setlength{\extrarowheight}{2pt}
                    \begin{tabular}{|l|l|}
                        \hline
                        \textbf{Command} & \textbf{Example} \\
                        \hline \hline
                        \texttt{text} & Normal text \\
                        \texttt{\textbackslash textnormal\{text\}} & \textnormal{Normal text} \\
                        \texttt{\textbackslash textbf\{text\}} & \textbf{Bold text} \\
                        \texttt{\textbackslash textit\{text\}} & \textit{Italic text} \\
                        \texttt{\textbackslash emph\{text\}} & \emph{Emphasized} \\
                        \texttt{\textbackslash underline\{text\}} & \underline{Underline} \\
                        \texttt{\textbackslash textsc\{text\}} & \textsc{Small Caps} \\
                        \texttt{\textbackslash texttt\{text\}} & \texttt{Typewriter family} \\
                        \texttt{\textbackslash textrm\{text\}} & \textrm{Roman family} \\
                        \texttt{\textbackslash textsf\{text\}} & \textsf{Sans serif family} \\
                        \hline
                    \end{tabular}
                \end{table}

            \subsection{ภาษาไทย}

                \begin{table}[h]
                    \centering
                    \setlength{\extrarowheight}{2pt}
                    \begin{tabular}{|l|l|}
                        \hline
                        \textbf{Command} & \textbf{Example} \\
                        \hline \hline
                        \texttt{text} & ข้อความ \\
                        \texttt{\textbackslash textnormal\{text\}} & \textnormal{ข้อความ} \\
                        \texttt{\textbackslash textbf\{text\}} & \textbf{ข้อความ} \\
                        \texttt{\textbackslash textit\{text\}} & \textit{ข้อความ} \\
                        \texttt{\textbackslash emph\{text\}} & \emph{ข้อความ} \\
                        \texttt{\textbackslash underline\{text\}} & \underline{ข้อความ} \\
                        \hline
                    \end{tabular}
                \end{table}

            \subsection{Sample Text}

                In Bangkok, \textbf{everywhere you look} is amazing food. To sound like a local, say \textit{bhed nid noi}
                (\textit{เผ็ดนิดหน่อย}), meaning \emph{a little bit spicy}.

                \vspace{0.5em}

                Navigate the city using \texttt{Google Maps} or the \textsf{BTS Skytrain} (\textnormal{รถไฟฟ้า}).
                If writing a formal journal, you might start with \textsc{Chapter One} in the standard \textrm{Roman font family}.
                Wherever you go, remember the most important phrase: \underline{Sawatdee krap} (\textnormal{สวัสดีครับ})!

    \pagebreak

        \section{Font Size}

            Font sizes are identified by special names, the actual size is not absolute but relative to the font size declared in the 
            \texttt{\textbackslash documentclass} statement.

            \subsection{English}
                \begin{table}[h]
                    \centering
                    \setlength{\extrarowheight}{8pt}
                    \begin{tabular}{|l|l|}
                        \hline
                        \textbf{Command} & \textbf{Example} \\
                        \hline \hline
                        \texttt{\{\textbackslash tiny text\}} & {\tiny Tiny} \\
                        \texttt{\{\textbackslash scriptsize text\}} & {\scriptsize Script size} \\
                        \texttt{\{\textbackslash footnotesize text\}} & {\footnotesize Footnote size} \\
                        \texttt{\{\textbackslash small text\}} & {\small Small} \\
                        \texttt{text} & Normal size \\
                        \texttt{\{\textbackslash normalsize text\}} & {\normalsize Normal size} \\
                        \texttt{\{\textbackslash large text\}} & {\large large} \\
                        \texttt{\{\textbackslash Large text\}} & {\Large Large} \\
                        \texttt{\{\textbackslash LARGE text\}} & {\LARGE LARGE} \\
                        \texttt{\{\textbackslash huge text\}} & {\huge huge} \\
                        \texttt{\{\textbackslash Huge text\}} & {\Huge Huge} \\
                        \hline
                    \end{tabular}
                \end{table}

            \subsection{ภาษาไทย}
                \begin{table}[h]
                    \centering
                    \setlength{\extrarowheight}{8pt}
                    \begin{tabular}{|l|l|}
                        \hline
                        \textbf{Command} & \textbf{Example} \\
                        \hline \hline
                        \texttt{\{\textbackslash tiny text\}} & {\tiny เล็กมากที่สุด} \\
                        \texttt{\{\textbackslash scriptsize text\}} & {\scriptsize เล็กมาก} \\
                        \texttt{\{\textbackslash footnotesize text\}} & {\footnotesize เล็กลง} \\
                        \texttt{\{\textbackslash small text\}} & {\small เล็ก} \\
                        \texttt{text} & ปกติ \\
                        \texttt{\{\textbackslash normalsize text\}} & {\normalsize ปกติ} \\
                        \texttt{\{\textbackslash large text\}} & {\large ใหญ่} \\
                        \texttt{\{\textbackslash Large text\}} & {\Large ใหญ่ขึ้น} \\
                        \hline
                    \end{tabular}
                \end{table}

    \pagebreak

                \begin{table}[h]
                    \centering
                    \setlength{\extrarowheight}{8pt}
                    \begin{tabular}{|l|l|}
                        \hline
                        \textbf{Command} & \textbf{Example} \\
                        \hline \hline
                        \texttt{\{\textbackslash LARGE text\}} & {\LARGE ใหญ่มาก} \\
                        \texttt{\{\textbackslash huge text\}} & {\huge ใหญ่มากมาก} \\
                        \texttt{\{\textbackslash Huge text\}} & {\Huge ใหญ่ที่สุด} \\
                        \hline
                    \end{tabular}
                \end{table}

\end{document}
